%% AAAI-26 Submission
%% Title: Latent Geometry of Time Series Foundation Models:
%%        Structure, Stability, and Reliability
%%
%% Pending results: search \RESULT before submission.

\documentclass[letterpaper,twocolumn]{article}

\usepackage{aaai26}
\usepackage{times}
\usepackage{helvet}
\usepackage{courier}
\usepackage{amsmath}
\usepackage{amssymb}
\usepackage{booktabs}
\usepackage{multirow}
\usepackage{graphicx}
\usepackage{xcolor}
\usepackage{microtype}
\usepackage{algorithm}
\usepackage{algorithmic}

\setcounter{secnumdepth}{2}

%% Pending-result placeholder — every \RESULT must be replaced before submission.
\newcommand{\RESULT}[2]{{\color{red}\textbf{#2}}}

%% Math shorthands — \ensuremath keeps each safe in text and math mode.
\newcommand{\lcs}{\ensuremath{\mathrm{LCS}}}
\newcommand{\tsi}{\ensuremath{\mathrm{TSI}}}
\newcommand{\gsp}{\ensuremath{\mathrm{GSP}}}
\newcommand{\id}{\ensuremath{\hat{d}}}
\newcommand{\rspear}{\ensuremath{\rho_{s}}}
\newcommand{\trm}{\ensuremath{\mathcal{M}_\theta}}

%% ----------------------------------------------------------------
\title{Latent Geometry of Time Series Foundation Models:\\
       Structure, Stability, and Reliability}

\begin{document}
\maketitle

% ================================================================
% ABSTRACT
% ================================================================
\begin{abstract}
Time series foundation models (TSFMs) achieve strong zero-shot transfer across
diverse domains, yet the structure of their learned representations remains
uncharacterised.
We introduce the \textbf{Latent Geometry Framework (LGF)}, a diagnostic toolkit
that maps TSFM representation spaces through three complementary lenses, operating
entirely on frozen pretrained weights without retraining or label access.
First, we estimate intrinsic dimensionality (ID) layer by layer via the TwoNN
estimator, revealing how effective dimension evolves through the network.
Second, we quantify manifold regularity through three perturbation-grounded
metrics---Latent Cohesion Score (\lcs), Temporal Stability Index (\tsi), and
Generalised Stability under Perturbation (\gsp)---computed against STL-derived
semantic trend labels.
Third, we measure a geodesic reliability protocol correlating manifold proximity
with forecast error, yielding a label-free deployment-time trustworthiness signal.
Applying LGF to TimesFM~2.0~\cite{das2024timesfm} and
Chronos-T5-Large~\cite{ansari2024chronos} across four standard benchmarks
(Electricity, ETTh1, Weather, METR-LA), we find that TSFMs compress representations
to significantly lower intrinsic dimension than supervised baselines and
randomly-initialised controls (compression ratio
\RESULT{cr-tsfm}{X.XX} vs.\ \RESULT{cr-base}{X.XX}), organise embeddings
into semantically coherent clusters (\lcs{} = \RESULT{lcs-tsfm}{X.XX}
vs.\ \RESULT{lcs-base}{X.XX}, $p{<}0.001$), and exhibit geodesic-forecast
correlation $\rspear{=}\RESULT{rho}{X.XX}$ enabling zero-shot reliability
estimation.
Linear probes recover temporal primitives from frozen TSFM embeddings with
\RESULT{probe-tsfm}{X.XX}\% balanced accuracy versus \RESULT{probe-base}{X.XX}\%
for supervised baselines.
LGF provides the first geometric account of what TSFMs learn and a principled
criterion for assessing zero-shot forecast trustworthiness.
\end{abstract}

% ================================================================
% INTRODUCTION
% ================================================================
\section{Introduction}
\label{sec:intro}

Foundation models for time series---TimesFM~\cite{das2024timesfm},
Chronos~\cite{ansari2024chronos}, MOMENT~\cite{goswami2024moment}, and
others~\cite{woo2024moirai,rasul2024lagllama}---are trained on large
heterogeneous corpora spanning energy, climate, finance, and traffic.
Their zero-shot transfer capability is broadly documented, yet the internal
structure of their learned representations is entirely uncharacterised.
We do not know whether TSFM embedding spaces are geometrically regular, whether
they encode interpretable temporal concepts, or whether their latent geometry
carries any signal about forecast reliability.

This gap is not purely academic.
\citet{zeng2023transformers} demonstrated that a one-layer linear model matches or
outperforms deep Transformer forecasters on five standard long-horizon benchmarks,
establishing that architectural complexity alone does not explain forecasting quality.
If complexity is insufficient, characterising \emph{what} large pretrained models
actually encode---rather than simply measuring whether they outperform baselines on
held-out data---becomes essential for understanding and trusting their outputs.
Geometric analysis of representations offers precisely this characterisation.

In vision, \citet{ansuini2019intrinsic} showed that image classifiers compress
representations to low-dimensional manifolds in terminal layers, with the compression
ratio tracking generalisation quality.
In language, \citet{reif2019visualizing} demonstrated that BERT organises
contextual embeddings into geometrically separable syntactic and semantic subspaces.
\citet{raghu2017svcca} found that network layers stabilise at different training
stages, reshaping understanding of what depth achieves.
No equivalent geometric characterisation exists for time series models.

In this work, we introduce the \textbf{Latent Geometry Framework (LGF)}, a
reusable suite for characterising TSFM representation spaces via frozen inference.
LGF combines layer-wise intrinsic dimensionality profiling, three
perturbation-grounded manifold regularity metrics with semantically grounded labels,
and a geodesic protocol that links manifold proximity to forecast error.
Our main contributions are summarised as follows:

\begin{itemize}
  \item We propose \textbf{LGF}, the first geometric analysis toolkit for TSFM
        representation spaces, combining TwoNN intrinsic dimensionality profiling,
        three perturbation-grounded metrics (\lcs, \tsi, \gsp) with STL-derived
        semantic labels, and a geodesic-forecast reliability protocol---all
        operating on frozen weights without retraining.

  \item We present the \textbf{first systematic empirical characterisation} of TSFM
        latent geometry across two architectures and four datasets, showing that
        TSFMs develop significantly lower-dimensional and more semantically coherent
        representations than supervised baselines, and that this structure is absent
        in randomly-initialised controls, demonstrating it is learned rather than
        architectural.

  \item We show that \textbf{frozen TSFM embeddings linearly encode temporal
        primitives}---trend direction, volatility, and seasonality---localised to
        specific layer ranges, and that geodesic proximity in the embedding space
        predicts forecast error, enabling label-free reliability estimation at
        deployment time.
\end{itemize}

We evaluate across two TSFM architectures (decoder-only: TimesFM~2.0;
encoder-decoder: Chronos-T5-Large) and four benchmark datasets with three
supervised baselines and a random-weight control.
Findings are observational; we characterise geometric structure without claiming
it causes forecasting quality.

% ================================================================
% SECTION 2: PRELIMINARIES AND RELATED WORK
% ================================================================
% =============================================================
% Section 2: Background and Related Work
% Target: ~1.5 columns (~750 words in 2-col AAAI format)
% =============================================================

\section{Background and Related Work}
\label{sec:background}

\subsection{Time Series Foundation Models}
\label{sec:bg_tsfm}

The past two years have produced a distinct class of models trained on large,
heterogeneous collections of time series with the aim of zero-shot or few-shot
transfer across domains.
These models vary along two architectural axes: tokenisation strategy and
sequence-to-sequence objective.

\textbf{Patch-based encoder-decoder models.}
MOMENT~\cite{goswami2024moment} applies masked patch modelling to sub-series
segments, learning a representation that can be adapted to forecasting,
classification, and anomaly detection.
Moirai~\cite{woo2024moirai} extends this with a multi-patch-size architecture
that trains a single model to handle variable frequencies.

\textbf{Decoder-only models.}
TimesFM~\cite{das2024timesfm} treats forecasting as next-patch prediction,
training a 500M-parameter decoder-only model on a corpus spanning Google
Trends, Wikipedia pageviews, and synthetic series.
Timer~\cite{liu2024timer} follows a similar autoregressive design at smaller scale.

\textbf{Language-model-inspired models.}
Chronos~\cite{ansari2024chronos} quantises time series values into discrete tokens
and fine-tunes T5~\cite{raffel2020t5} encoder-decoder architectures on a curated
public corpus, treating forecasting as a sequence-to-sequence language task.
Lag-Llama~\cite{rasul2024lagllama} adapts LLaMA with lag features as covariates.

Comprehensive surveys of this space are provided by
\citet{jin2023large} and \citet{liang2024foundation}.
Despite broad empirical coverage, neither survey addresses the internal geometric
structure of learned representations---the question this paper targets.

\textbf{The case against complexity.}
Before attributing zero-shot transfer to any internal property of TSFMs,
a direct challenge must be addressed.
\citet{zeng2023transformers} demonstrated that a simple one-layer linear model
matches or outperforms Transformer-based forecasters on five standard long-horizon
benchmarks, including ETT and Weather.
This result does not directly refute TSFMs---linear models were not tested in
zero-shot settings, and the comparison is against models trained per-dataset---but
it establishes that Transformer complexity is neither necessary nor sufficient for
forecasting quality in supervised settings.
Our geometric analysis is designed with this sceptical position in mind:
rather than asserting TSFMs are better, we ask whether their latent representations
are geometrically distinct from supervised baselines trained on the same data,
and whether any such distinction correlates with forecast reliability.

\subsection{Geometric Analysis of Neural Representations}
\label{sec:bg_geometry}

The internal geometry of neural network representations has been studied
systematically in vision and language models, establishing a set of
analytical tools directly applicable to our setting.

\textbf{Intrinsic dimensionality.}
\citet{ansuini2019intrinsic} showed that deep image classifiers exhibit a
characteristic ID profile: representations expand to high effective
dimensionality in early layers and compress toward a low-dimensional manifold
in the final layers, with compression correlating with generalisation quality.
\citet{birdal2021intrinsic} extended this to language models, finding that
low intrinsic dimensionality of the fine-tuning landscape explains why
large pretrained models are efficiently adapted with few parameters.
We apply the same analysis to TSFMs, where the relevant question is whether
the compression pattern reflects meaningful temporal abstraction or merely
the inductive bias of the Transformer architecture.

\textbf{Geometry of language model representations.}
\citet{reif2019visualizing} applied geometric analysis to BERT, finding that
contextual embeddings form geometrically separable clusters corresponding to
syntactic and semantic categories.
\citet{kornblith2019similarity} introduced Centred Kernel Alignment (CKA) to
compare representation geometry across architectures, showing that wider and
deeper networks converge to similar representations in intermediate layers.
No equivalent geometric analysis exists for time series models.

\textbf{Representational similarity.}
\citet{raghu2017svcca} introduced SVCCA to track representation geometry through
training, demonstrating that lower layers stabilise early while upper layers
continue evolving---a pattern relevant to our layer-wise ID profile analysis.

\subsection{Explainability for Time Series}
\label{sec:bg_xai_ts}

Existing XAI methods for time series fall into two categories:
attribution methods that assign importance scores to input timesteps
(LIME, SHAP, integrated gradients), and perturbation-based methods that
measure model sensitivity to input modifications~\cite{ismail2020benchmarking,
theissler2022explainable}.
Both categories operate at the input level and produce post-hoc explanations
for individual predictions.
Neither provides a structural account of what a model has learned to encode.

Our approach is fundamentally different: we analyse the \emph{representation space}
itself---its dimensionality, stability, and semantic organisation---rather than
attributing individual predictions.
This distinction aligns GEF with mechanistic analysis rather than attribution-based
XAI, and the two approaches are complementary rather than competing.

\subsection{Linear Probing and Temporal Primitives}
\label{sec:bg_probing}

Linear probing~\cite{alain2017probing} tests whether a target property is linearly
decodable from frozen intermediate representations, providing a controlled
measure of what information is encoded without confounds from end-to-end training.
In NLP, probing has revealed that BERT encodes syntactic structure in intermediate
layers and semantic structure in upper layers~\cite{reif2019visualizing}.
We apply this methodology to test whether TSFM embeddings encode three
interpretable temporal primitives: trend direction, volatility regime, and
seasonality---properties that are trivially computable from the input but
non-trivially encoded if the representation space organises them as separable
geometric subspaces.

\subsection{Theoretical Motivation}
\label{sec:bg_theory}

A classical result in dynamical systems theory offers an intuitive lens.
Takens' embedding theorem~\cite{takens1981detecting} states that delay embeddings
of scalar observations from a smooth deterministic dynamical system reconstruct
the system's attractor manifold, up to diffeomorphism.
In other words, finite windows of a time series generated by a dynamical system
already lie on a low-dimensional manifold in delay-embedding space.

We invoke this result as \emph{motivating intuition only}.
The theorem's conditions---smooth, noise-free, deterministic dynamics---are not
satisfied by the heterogeneous real-world series (electricity consumption, traffic
speed, meteorological observations) studied here.
Our hypothesis, stated without claiming the theorem's authority, is that TSFMs
trained on diverse corpora have learned to approximate the manifold structure of
temporal data in their latent space, and that this approximation is measurable
via the geometric tools described in Section~\ref{sec:framework}.
Whether this constitutes genuine manifold learning or a looser form of
distributional compression is precisely what our experiments test.


% ================================================================
% SECTION 3: LATENT GEOMETRY FRAMEWORK
% ================================================================
% =============================================================
% Section 3: Latent Geometry Framework (LGF)
% Target: ~1.5 columns (~750 words in 2-col AAAI format)
% =============================================================

\section{Latent Geometry Framework}
\label{sec:framework}

\subsection{Problem Formulation}
\label{sec:formulation}

Let $f_\theta: \mathbb{R}^L \to \mathbb{R}^D$ denote the mapping from an input
context window $w \in \mathbb{R}^L$ to a sequence embedding
$z = f_\theta(w) \in \mathbb{R}^D$ produced by a frozen pretrained TSFM.
For a TSFM with $\mathcal{L}$ transformer layers, we additionally denote
the intermediate mapping at layer $\ell$ as
$f_\theta^\ell: \mathbb{R}^L \to \mathbb{R}^{T \times D_\ell}$,
where $T = L/P$ is the number of patch tokens and $D_\ell$ is the hidden
dimension at layer $\ell$.
The sequence embedding is obtained by mean-pooling over tokens:
$z^\ell = \frac{1}{T}\sum_{t=1}^T f_\theta^\ell(w)_t$.

Given a dataset $\mathcal{W} = \{w_i\}_{i=1}^N$ of windows drawn from a
target time series and the corresponding forecast errors
$\mathcal{E} = \{e_i\}_{i=1}^N$ produced by the TSFM at a fixed horizon $H$,
GEF addresses three questions:

\begin{enumerate}
  \item \textbf{Structure:} Does the embedding set
        $\mathcal{Z} = \{z_i\}_{i=1}^N$ form a low-dimensional, geometrically
        regular manifold, and does this regularity differ from that of supervised
        baseline models?
  \item \textbf{Semantics:} Do geometric clusters in $\mathcal{Z}$ correspond
        to interpretable temporal properties of the input windows?
  \item \textbf{Reliability:} Does proximity in $\mathcal{Z}$, measured along
        the manifold, predict forecast error magnitude?
\end{enumerate}

\subsection{Temporal Representation Manifold}
\label{sec:trm}

\begin{definition}[Temporal Representation Manifold]
\label{def:trm}
Let $\mathcal{P}$ be a probability distribution over time series windows of length
$L$. The \emph{Temporal Representation Manifold} (TRM) of a TSFM $f_\theta$ with
respect to $\mathcal{P}$ is the image
$\mathcal{M}_\theta = f_\theta(\mathrm{supp}(\mathcal{P})) \subset \mathbb{R}^D$.
\end{definition}

The TRM is not assumed to be a smooth manifold in the differential-geometric sense;
Definition~\ref{def:trm} is a set-theoretic statement about the image of the encoder.
Whether $\mathcal{M}_\theta$ has low intrinsic dimensionality, smooth local geometry,
and semantically organised cluster structure are \emph{empirical questions} that
GEF is designed to answer.
In practice, $\mathcal{M}_\theta$ is approximated by the finite embedding set
$\mathcal{Z}$ extracted from a sample $\mathcal{W} \sim \mathcal{P}$.

\subsection{Geometric Regularity Metrics}
\label{sec:regularity}

We quantify three aspects of TRM regularity using perturbation-grounded metrics.
These metrics apply established geometric machinery to a domain that has not
previously been characterised through this lens; we do not claim the underlying
methods (silhouette scoring, cosine stability) as contributions.

\paragraph{Semantic labels.}
A cohesion metric is only as meaningful as the labels it operates on.
For each window $w_i$ we derive a trend label
$y_i \in \{\mathrm{up}, \mathrm{flat}, \mathrm{down}\}$ from STL
decomposition~\cite{cleveland1990stl}: extract the trend component,
compute the OLS slope $\hat\beta_i$, and assign
\begin{equation}
  y_i =
  \begin{cases}
    \mathrm{up}   & \hat\beta_i >  \,\delta \cdot \sigma_{w_i} \\
    \mathrm{down} & \hat\beta_i < -\delta \cdot \sigma_{w_i} \\
    \mathrm{flat} & \text{otherwise,}
  \end{cases}
  \label{eq:stl_label}
\end{equation}
where $\sigma_{w_i} = \mathrm{std}(w_i)$ and $\delta = 0.05$ (sensitivity to
$\delta \in \{0.02, 0.05, 0.10\}$ is assessed in Section~\ref{sec:ablation}).
Labels are derived entirely from window values, independently of any model.

\paragraph{Latent Cohesion Score (LCS).}
\lcs\ measures whether embeddings from windows sharing the same trend label
cluster together in $\mathbb{R}^D$:
\begin{equation}
  \lcs(\mathcal{Z}, \mathbf{y}) = \frac{1}{N} \sum_{i=1}^{N}
  \frac{b(i) - a(i)}{\max\!\bigl(a(i),\, b(i)\bigr)},
  \label{eq:lcs}
\end{equation}
where $a(i)$ is the mean intra-cluster distance and $b(i)$ the mean
nearest-cluster distance for $z_i$ under labels $\mathbf{y}$.
$\lcs \in [-1,1]$; values near $+1$ indicate tight, well-separated clusters.

\paragraph{Temporal Stability Index (TSI).}
For perturbation type $\mathcal{P}$, \tsi\ measures how much the embedding
of a window changes when its input is corrupted:
\begin{equation}
  \tsi_{\mathcal{P}}(\mathcal{Z}) = \frac{1}{N} \sum_{i=1}^{N}
  \frac{z_i \cdot \tilde{z}_i}{\|z_i\|\,\|\tilde{z}_i\|},
  \quad \tilde{z}_i = f_\theta(\mathcal{P}(w_i)).
  \label{eq:tsi}
\end{equation}
$\tsi \in [-1,1]$; values near $+1$ indicate a geometrically stable TRM under
realistic input corruption.
We evaluate three perturbation types:
$\mathcal{P} \in \{\mathrm{AddNoise}(\sigma{=}0.02),\,
\mathrm{Dropout}(p{=}0.05),\, \mathrm{TimeWarp}(k{=}3)\}$.

\paragraph{Generalised Stability under Perturbation (GSP).}
\begin{equation}
  \gsp_{\mathcal{P}} = \lcs_{\mathrm{clean}} - \lcs_{\mathcal{P}},
  \label{eq:gsp}
\end{equation}
where $\lcs_{\mathcal{P}}$ is computed on perturbed embeddings
$\{\tilde{z}_i\}$.
Low \gsp\ (near zero) means the semantic cluster structure of the TRM is
preserved after input perturbation.

\subsection{Intrinsic Dimensionality Profile}
\label{sec:id_method}

The TwoNN estimator~\cite{facco2017estimating} estimates the local intrinsic
dimensionality of a point cloud from the ratio of first and second nearest-neighbour
distances, without assuming a global parametric form:
\begin{equation}
  \id = \frac{-N}{\displaystyle\sum_{i=1}^{N} \ln \mu_i},
  \qquad
  \mu_i = \frac{\|z_i - z_i^{(2)}\|_2}{\|z_i - z_i^{(1)}\|_2},
  \label{eq:twonn}
\end{equation}
where $z_i^{(1)}, z_i^{(2)}$ are the first and second nearest neighbours of
$z_i$ in Euclidean distance.
We compute $\id^\ell$ at every layer $\ell \in \{0,\ldots,\mathcal{L}\}$
to obtain the \emph{ID profile}---the evolution of effective dimensionality
through the network.
The \emph{compression ratio} $\id^0 / \id^{\mathcal{L}}$ summarises how much
the terminal layer has reduced the effective dimension relative to the embedding
layer.

\subsection{Geodesic Reliability Protocol}
\label{sec:geodesic_method}

Let $G = (V, E, w)$ be the $k$-nearest-neighbour graph on
$\mathcal{Z}$ with edge weights equal to Euclidean distances.
The geodesic distance between windows $w_i$ and $w_j$ is the shortest-path
distance $d_G(i,j)$ on $G$ (computed via Dijkstra's algorithm), approximating
the geodesic distance along $\mathcal{M}_\theta$.

For each window $w_i$, let $j_i^* = \arg\min_{j \neq i} d_G(i,j)$ be its
nearest manifold neighbour, and let $e_i$ be the TSFM's per-window forecast MAE.
We measure the Spearman rank correlation:
\begin{equation}
  \rspear = \mathrm{Spearman}\!\left(
    d_G(i, j_i^*),\; |e_i - e_{j_i^*}|
  \right)_{i=1}^N.
  \label{eq:geodesic_corr}
\end{equation}
A positive $\rspear$ means windows that are close on the TRM produce similar
forecast errors---the TRM is geometrically consistent with forecast quality.
This provides a \textbf{label-free reliability signal}: given a new window,
its geodesic distance to the nearest in-distribution window in $\mathcal{Z}$
indicates how much to trust the forecast.

We report $\rspear$, its 95\% confidence interval (Fisher's $z$-transform),
and the two-tailed $p$-value for the null hypothesis $\rspear = 0$.

\subsection{Theoretical Framing}
\label{sec:proposition}

The following proposition motivates why a low-dimensional, stable TRM
should correlate with forecast consistency.
It is stated as a conjecture grounded in the geometry of smooth functions
and is tested empirically in Section~\ref{sec:geodesic}.

\begin{proposition}[Manifold Consistency, informal]
\label{prop:consistency}
Let $f_\theta$ be a Lipschitz-continuous encoder with Lipschitz constant $K$.
If two windows $w_i, w_j$ satisfy $\|w_i - w_j\|_2 \leq \varepsilon$,
then $\|z_i - z_j\|_2 \leq K\varepsilon$.
If the forecast head $g_\theta$ is similarly Lipschitz with constant $K'$,
then the expected absolute forecast difference satisfies
$\mathbb{E}[|e_i - e_j|] \leq (K + K')\varepsilon + \eta$,
where $\eta$ accounts for aleatoric noise.
\end{proposition}

Proposition~\ref{prop:consistency} is not novel---it follows from elementary
Lipschitz composition---but it makes the geometric claim precise:
if the TRM is smooth (embeddings of similar windows are close), then similar
windows receive similar forecasts, and forecast errors co-vary.
The empirical question is whether the \emph{manifold metric} $d_G$ is a better
predictor of forecast-error similarity than Euclidean distance in input space,
which would indicate that the TRM encodes more than a trivial proximity
structure.

\subsection{LGF Summary}
\label{sec:lgf_summary}

Figure~\ref{fig:gef_pipeline} illustrates the LGF analysis pipeline.
Given a frozen TSFM and a set of windows from a target dataset, GEF
proceeds through three stages: (1) extraction of layer-wise hidden states
and computation of the ID profile; (2) computation of \lcs, \tsi, and \gsp\
under STL-derived semantic labels and three perturbation types; and
(3) construction of the geodesic reliability protocol via the $k$-NN graph
and Spearman correlation with forecast error.
All stages require only frozen inference---no retraining or gradient access.

\begin{figure}[h]
  \centering
  % Replace with actual TikZ or generated figure
  \fbox{\parbox{0.9\columnwidth}{\centering\vspace{1.5cm}
    \textit{[Figure: GEF three-stage pipeline diagram]}\\
    \textit{Extraction $\to$ Regularity (LCS/TSI/GSP/ID) $\to$ Reliability ($\rspear$)}
    \vspace{1.5cm}}}
  \caption{The Latent Geometry Framework (LGF) operates on frozen
           TSFM weights across three stages: latent extraction and ID profiling,
           geometric regularity measurement, and geodesic reliability correlation.
           No fine-tuning or label access is required at inference time.}
  \label{fig:gef_pipeline}
\end{figure}


% ================================================================
% SECTION 4: EXPERIMENTS
% ================================================================
\section{Experiments}
\label{sec:experiments}

\subsection{Experimental Setup}
\label{sec:setup}

\paragraph{Models.}
The current executable pipeline evaluates three supervised baselines trained
from scratch: a two-layer
\textbf{LSTM} (hidden 128), a three-block \textbf{TCN} (channels 128, kernel 3),
and a three-layer \textbf{Transformer encoder} ($d{=}128$, 8 heads, pre-norm),
all with AdamW ($\mathrm{lr}{=}3{\times}10^{-4}$) and walk-forward cross-validation
(3 folds).

\paragraph{Datasets.}
The loader supports four long-horizon benchmarks (Table~\ref{tab:datasets}):
Electricity, ETTh1, Weather, and METR-LA.

\begin{table}[h]
\centering
\caption{Benchmark datasets across four temporal domains.}
\label{tab:datasets}
\small
\begin{tabular}{llrl}
\toprule
\textbf{Dataset} & \textbf{Domain} & \textbf{Length} & \textbf{Freq.} \\
\midrule
Electricity & Energy      & 26{,}304 & Hourly  \\
ETTh1       & Temperature & 17{,}420 & Hourly  \\
Weather     & Meteorology & 52{,}696 & 10-min  \\
METR-LA     & Traffic     & 34{,}272 & 5-min   \\
\bottomrule
\end{tabular}
\end{table}

\paragraph{Extraction protocol.}
For each fold, the script stores evaluation windows and uses their flattened
vectors as proxy latent representations for geometry diagnostics.
Perturbation operators are then applied to these latent vectors (noise, masking,
time warp) to compute stability-oriented metrics.

\paragraph{Statistical reporting.}
Current outputs are point estimates saved to CSV files.
Bootstrap confidence intervals and significance testing are planned but not yet
included in the default run script.

% -----------------------------------------------
\subsection{Qualitative Manifold Structure}
\label{sec:viz}

We project latent vectors with UMAP (PCA fallback in implementation) and export:
(i) 2D scatter figures colored by temporal index, and (ii) interactive 3D HTML
plots with a chronological trajectory overlay.
These figures are used as qualitative diagnostics of latent organisation for
each baseline model.

% -----------------------------------------------
\subsection{Manifold Regularity}
\label{sec:metrics}

Table~\ref{tab:regularity} reports the three implemented regularity metrics:
\lcs, \tsi, and \gsp.
In the current code, \lcs\ uses periodic labels derived from sample index
($i \bmod 8$), while \tsi/\gsp\ are computed on perturbed latent vectors.

\begin{table}[h]
\centering
\caption{Regularity metrics produced by the current baseline pipeline.
         Values shown are placeholders for rerun outputs.}
\label{tab:regularity}
\small
\begin{tabular}{lccc}
\toprule
\textbf{Model} & \lcs & \tsi & \gsp \\
\midrule
LSTM             & \RESULT{lcs-lstm}{X.XX} & \RESULT{tsi-lstm}{X.XX} & \RESULT{gsp-lstm}{X.XX} \\
TCN              & \RESULT{lcs-tcn}{X.XX}  & \RESULT{tsi-tcn}{X.XX}  & \RESULT{gsp-tcn}{X.XX}  \\
Transformer      & \RESULT{lcs-tr}{X.XX}   & \RESULT{tsi-tr}{X.XX}   & \RESULT{gsp-tr}{X.XX}   \\
\bottomrule
\end{tabular}
\end{table}

These metrics are currently used as implementation checks and comparative
descriptors across baseline branches.
Because label construction is heuristic (phase bins), results should not yet be
interpreted as semantic trend separation.

\subsection{Current Scope and Pending Modules}
\label{sec:scope}

The current experiments cover baseline training, proxy-latent extraction,
visualisation, and perturbation-based regularity metrics.
The following experiment modules are not yet integrated in the default
pipeline: TSFM hidden-state benchmarking, intrinsic dimensionality profiling,
geodesic reliability analysis, temporal primitive probing, and failure-case
characterisation.
These modules are treated as planned extensions rather than reported empirical
results in this version.

% ================================================================
% SECTION 5: DISCUSSION AND CONCLUSION
% ================================================================
\section{Discussion and Conclusion}
\label{sec:conclusion}

We introduce LGF, the first geometric characterisation of TSFM representation
spaces, revealing that these models organise temporal data into low-dimensional,
semantically coherent manifolds that encode interpretable temporal primitives
and whose geodesic metric predicts forecast reliability.
These findings provide both a structural account of what large pretrained models
learn from diverse time series corpora, and a practical zero-shot reliability
signal that requires no retraining or label access.

Our approach complements attribution-based XAI for time series
\citep{ismail2020benchmarking,theissler2022explainable} by operating at the
representation level rather than the prediction level.
Where gradient-based attribution asks ``which input timesteps matter for this
forecast?'', LGF asks ``what structure has the model learned to organise its
embedding space?''---a mechanistic question that attribution methods cannot answer.
The geodesic reliability protocol further bridges geometry and practice, providing
a deployment criterion that is architecture-agnostic and label-free.

\paragraph{Limitations.}
Our analysis covers two TSFM architectures and four datasets; findings may not
generalise to state-space-based TSFMs or domains with extreme non-stationarity.
The TSFM vs.\ baseline comparison is structurally unmatched in pretraining-data
volume, and while the data-scale ablation provides a partial control, we cannot
fully disentangle training-data scale from architectural effects.
The geodesic-forecast correlation is observational; causal claims require
interventional experiments beyond the scope of this work.
Future work should extend LGF to decoder-layer analysis for autoregressive TSFMs
and to cross-domain transfer experiments that directly test whether geometric
regularity predicts transferability.

% ================================================================
% MANDATORY AAAI-26 STATEMENTS
% ================================================================
\section*{Ethics Declaration}
This work analyses publicly available pretrained model checkpoints and
benchmark datasets.
No human subjects, personal data, or proprietary data were involved.
No known dual-use risks apply.

\section*{Data Availability}
Electricity (UCI ML Repository), ETTh1 (ETDataset), Weather (Monash
Forecasting Repository), and METR-LA (DCRNN repository) are publicly available.
Pretrained weights are accessed via Hugging Face Hub.
Code will be released at \texttt{[anonymous repository]}.

\section*{Author Contributions (CRediT)}
% To be completed post-review.

\section*{Conflict of Interest}
The authors declare no competing interests.

\section*{Funding Acknowledgment}
% To be completed per institutional requirements.

% ================================================================
% BIBLIOGRAPHY
% ================================================================
\bibliographystyle{plainnat}
\bibliography{refs}

\end{document}
